\VerdierPhysicalPageAnchor{0038}%
\VerdierSemanticAnchor{ch1-1-2-4-isomorphisms-cont}
et d'isomorphismes :

\[
  m(g)\colon F_eT(g)\longrightarrow T'(g)F_e,
  \qquad g\in G,
\]
tels que les diagrammes (1.2.3.2) soient commutatifs.

\VerdierSemanticAnchor{ch1-1-2-4-characterization}
Un $G^\circ$-foncteur est donc un foncteur
“commutant à isomorphisme près aux opérations de $G^\circ$”. On
utilisera pour le désigner la notation $(F_e,m)$.

\VerdierSemanticAnchor{ch1-1-2-4-strict}
Lorsque pour tout $g\in G$ les isomorphismes $m(g)$ sont des
isomorphismes identiques, ce qui implique en particulier que pour tout
$g\in G$ on a les égalités $F_eT(g)=T'(g)F_e$, on dit que le
$G^\circ$-foncteur est \emph{strict}.

\VerdierSemanticAnchor{ch1-1-2-5}
\noindent\textbf{1.2.5.}\quad
Soient $\mathcal C$, $\mathcal C'$ et $\mathcal C''$ trois catégories
graduées et
\[
  \mathcal C\xrightarrow{F}\mathcal C'
  \xrightarrow{F'}\mathcal C'',\qquad F''=F'F,
\]
trois foncteurs gradués. Choisissons sur $\mathcal C$, $\mathcal C'$ et
$\mathcal C''$ des clivages normalisés. On en déduit, d'après ce qui
précède, trois $G^\circ$-catégories
$(\mathcal C_e,G^\circ,T,c)$,
$(\mathcal C'_e,G^\circ,T',c')$ et
$(\mathcal C''_e,G^\circ,T'',c'')$, et trois $G^\circ$-foncteurs
$(F_e,m)$, $(F'_e,m')$ et $(F''_e,m'')$. On vérifie immédiatement qu'on a
les relations :

\VerdierSemanticAnchor{ch1-1-2-5-eq1}
\begin{equation}
  \begin{cases}
    F''_e=F'_eF_e,\\[0.5em]
    m''(g)=\bigl(m'(g)\star F_e\bigr)
      \bigl(F'_e\star m(g)\bigr),\qquad g\in G.
  \end{cases}
  \tag{1.2.5.1}
\end{equation}

\VerdierSemanticAnchor{ch1-1-2-6}
\noindent\textbf{Définition 1.2.6.}\quad
Soient $(\mathcal C_e,G^\circ,T,c)$,
$(\mathcal C'_e,G^\circ,T',c')$ et
$(\mathcal C''_e,G^\circ,T'',c'')$ trois $G^\circ$-catégories et :

\[
\begin{tikzcd}[column sep=large]
(\mathcal C_e,G^\circ,T,c)
  \arrow[r,"{(F_e,m)}"]
&
(\mathcal C'_e,G^\circ,T',c')
  \arrow[r,"{(F_e^{\prime},m^{\prime})}"]
&
(\mathcal C''_e,G^\circ,T'',c'')
\end{tikzcd}
\]
deux $G^\circ$-foncteurs.

\VerdierSemanticAnchor{ch1-1-2-6-composite}
Les égalités (1.2.5.1) définissent un $G^\circ$-foncteur qui est appelé
le \emph{foncteur composé} des $G^\circ$-foncteurs $(F_e,m)$ et
$(F'_e,m')$.

\VerdierSemanticAnchor{ch1-1-2-7}
\noindent\textbf{1.2.7.}\quad
Soient $F$ et $F'$ deux foncteurs gradués entre deux catégories
$\mathcal C$ et $\mathcal C'$ graduées de type $G$, et soit
$u\colon F\to F'$ un morphisme de foncteurs de degré $e$. Choisissons
sur $\mathcal C$ et $\mathcal C'$ des clivages normalisés. On en déduit
deux $G^\circ$-catégories $(\mathcal C_e,G^\circ,T,c)$ et
$(\mathcal C'_e,G^\circ,T',c')$, et deux $G^\circ$-foncteurs
$(F_e,m)$, $(F'_e,m')$. Le morphisme $u$ détermine alors un morphisme de
foncteurs :

\VerdierSemanticAnchor{ch1-1-2-7-eq1}
\begin{equation}
  u_e\colon F_e\longrightarrow F'_e
  \tag{1.2.7.1}
\end{equation}
tel que le diagramme ci-après soit commutatif :

\VerdierSemanticAnchor{ch1-1-2-7-eq2}
\begin{equation}
\begin{tikzcd}[column sep=huge,row sep=large]
F_eT(g)
  \arrow[r,"u_e\star T(g)"]
  \arrow[d,"m(g)",swap]
&
F_e^{\prime}T(g)
  \arrow[d,"m^{\prime}(g)"]
\\
T^{\prime}(g)F_e
  \arrow[r,"T^{\prime}(g)\star u_e",swap]
&
T^{\prime}(g)F_e^{\prime}\mathpunct{.}
\end{tikzcd}
\tag{1.2.7.2}
\end{equation}
